\section{Strategy Phase}
\subsection{Arrival from Off-Table}

Some units have abilities that allow them to remain undeployed and enter the battlefield at this stage, most notably units with the \textit{Deep Strike} ability.

Each player rolls 1D6. The player with the highest result may place one of their detachments arriving from off-table, if they wish. Their opponent may then place one of theirs, and the players continue alternating in this manner.

A player may decide not to bring on any more detachments. Their opponent may then finish bringing on any remaining detachments they wish to deploy.

\subsection{Start-of-Turn Effects}

Some effects or special abilities are resolved at the very beginning of the turn. These are described in the relevant Codexes. For example, some armies have cards or effects that must be assigned; these are played at this stage.

\subsection{Assigning Orders}

NetEpic uses counters to indicate the orders assigned to detachments. These orders represent the actions the detachments may perform during the turn.

Each player selects an order for every detachment in their army that does not already have a Fall Back order. To do so, place an order counter face down beside each detachment, including off-table detachments that may enter during the turn. Some orders allow a choice between multiple actions, which is made when the detachment is activated.

If you forget to assign an order to a detachment, or if it has no order for another reason, it is considered to be \textit{without orders}.

There are four orders that may be assigned to a detachment:

\subsubsection{Advance Order}

These detachments move cautiously while retaining the ability to fire.

\begin{description}
    \item[Advance:] The detachment may move up to its Movement characteristic and fire normally during the Combat Phase. Once the movement has been completed, leave the revealed Advance order beside the detachment.
\end{description}

\subsubsection{Charge Order}

These detachments sacrifice their firepower in exchange for greater speed and the ability to engage the enemy in an assault. A Charge order allows one of two different actions:

\begin{description}
    \item[Charge:] The detachment may move up to twice its Movement characteristic and may engage enemy units in an assault. Once the action has been completed, remove the Charge order.

    \item[Countercharge:] The detachment may make a normal move and retains its revealed Charge order beside it. This allows the detachment to make a second move as a reaction, provided that this move enables it to engage an enemy detachment that ends its movement within range.
\end{description}

\subsubsection{Forced March Order}

These detachments focus on mobility and redeployment across the battlefield. A Forced March order allows one of two different actions:

\begin{description}
    \item[Rapid Advance:] The detachment may move up to twice its Movement characteristic and make Snap Fire during the Combat Phase. Once the movement has been completed, leave the revealed Forced March order beside the detachment.

    \item[Redeployment:] The detachment may make a triple move but cannot end its movement within an enemy interdiction zone. It may perform no further actions during the turn. If the detachment is a transport, the transported troops automatically receive a Redeployment order, replacing their current order. Once the movement has been completed, remove the order.
\end{description}

\subsubsection{First Fire Order}

These detachments sacrifice movement to gain an advantage when firing. Detachments with a First Fire order cannot move or change their facing. In return, they may reroll results of 1 on their ranged attacks' To-Hit rolls.

A First Fire order allows one of two different actions:

\begin{description}
    \item[First Fire:] First Fire allows the detachment to perform special attacks, such as indirect artillery fire and surprise attacks.

    \item[Overwatch Fire:] This action requires an unrevealed First Fire order and allows the detachment to react to enemy movement. The detachment may fire during the Movement Phase with a $-1$ modifier to its To-Hit rolls. See the Movement Phase section for further details.
\end{description}

\subsubsection{Fall Back Order / Without Orders}

Detachments that are falling back or have no orders may perform one of two different actions:

\begin{description}
    \item[Hold:] The detachment cannot move but may make Snap Fire during the Combat Phase. It receives a revealed Forced March order.

    \item[Regroup:] The detachment may make a normal move but cannot engage in an assault. Once the movement has been completed, remove its order counter, if it has one.
\end{description}

\begin{table}[ht]
    \centering

    \resizebox{\textwidth}{!}{%
        \begin{tblr}{
            colspec={
                Q[c,m,2]
                Q[c,m,2]
                Q[c,m,3.7]
                Q[c,m,3.4]
            },
            cells={
                halign=c,
                valign=m
            },
            hlines,
            vlines,
            row{1}={
                bg=black,
                fg=white,
                font=\bfseries,
                ht=10mm
            },
            row{2}={bg=green!40},
            row{3,4}={bg=cyan!55},
            row{5,6}={bg=violet!55},
            row{7,8}={bg=red!45},
            row{9,10}={bg=yellow!45}
        }
            Order Counter &
            Action &
            Movement Phase &
            Shooting during the Combat Phase
            \\

            \textit{\textbf{Advance}} &
            \textbf{Advance} &
            Normal move &
            Shooting
            \\

            \SetCell[r=2]{c,m}
                \textit{\textbf{Charge}} &
            \textbf{Charge} &
            \SetCell{c,m}{Double move allowing the detachment\\to engage in an assault} &
            --
            \\

            &
            \textbf{Countercharge} &
            \SetCell{c,m}{Normal move, followed by a possible\\Countercharge reaction} &
            --
            \\

            \SetCell[r=2]{c,m}
                \textit{\textbf{Forced March}} &
            \textbf{Rapid Advance} &
            Double move &
            Snap Fire
            \\

            &
            \textbf{Redeployment} &
            \SetCell{c,m}{Triple move, ending outside\\enemy interdiction zones} &
            --
            \\

            \SetCell[r=2]{c,m}
                \textit{\textbf{First Fire}} &
            \textbf{First Fire} &
            -- &
            First Fire
            \\

            &
            \textbf{Overwatch Fire} &
            Overwatch Fire as a reaction &
            --
            \\

            \SetCell[r=2]{c,m}
                \SetCell{font=\bfseries\itshape}
                {Fall Back\\Without Orders} &
            \textbf{Hold} &
            -- &
            Snap Fire
            \\

            &
            \textbf{Regroup} &
            Normal move &
            --
        \end{tblr}%
    }
\end{table}

\subsection{Initiative}

Each player rolls 1D6. The player with the highest result decides which player gains the initiative.

If the result is tied, the player who did not have the initiative during the previous turn gains it for the current turn. If the result is tied during the first turn, reroll the dice.